\documentclass[aspectratio=169]{beamer}
% Shared Beamer configuration for MUS2640 lecture decks (included after \documentclass).
\usepackage[utf8]{inputenc}
\usepackage[T1]{fontenc}
\usepackage{lmodern}
\usepackage{graphicx}
\usepackage{booktabs}

\usetheme{Madrid}
\usecolortheme{dolphin}
\usefonttheme{professionalfonts}

\setbeamercovered{transparent}
\setbeamertemplate{navigation symbols}{}
\setbeamertemplate{footline}[frame number]

\newcommand{\coursename}{MUS2640 · Sensing Sound and Music}
\newcommand{\instituteuo}{University of Oslo · Department of Musicology / RITMO}


\title{Week 11 · Vision}
\subtitle{\coursename}
\institute{\instituteuo}
\date{}

\begin{document}

\begin{frame}
  \titlepage
\end{frame}

\begin{frame}{Aims}
  \begin{itemize}
    \item Distinguish \textbf{audio/video} (signals) from \textbf{auditory/visual} (perception).
    \item Audiovisual integration — timing windows; McGurk-style lessons for music.
    \item Gaze, attention, and eye-tracking — what we measure and why it matters musically.
  \end{itemize}
\end{frame}

\begin{frame}{Roadmap}
  \begin{enumerate}
    \item Why vision changes heard timing and articulation in performance.
    \item Eye anatomy/optics sketch — sufficient for interpreting gaze papers.
    \item Attention: selective listening vs looking in multimodal scenes.
    \item Touch/vibration tie-in — multimodal beyond AV (bridge to earlier weeks).
  \end{enumerate}
\end{frame}

\begin{frame}{Design implications}
  \begin{itemize}
    \item Camera angle and lighting shape perceived musical expression.
    \item Streaming and VR extend ``presence'' — ethical and perceptual questions.
  \end{itemize}
\end{frame}

\begin{frame}{In-class}
  \begin{itemize}
    \item Short clip: note gaze targets vs listening goals (solo vs ensemble).
    \item Discuss when vision helps vs distracts from listening.
  \end{itemize}
\end{frame}

\begin{frame}{Outside class}
  \begin{itemize}
    \item Reading + short analysis of a performance clip (per weekly matrix).
  \end{itemize}
\end{frame}

\begin{frame}{Summary}
  \begin{itemize}
    \item Music perception is routinely \textbf{audiovisual}; methods must reflect that.
    \item Next: \textbf{Machine listening} — features, MIR, ML, human comparison, ethics.
  \end{itemize}
\end{frame}

\end{document}
